W kolejnej części pracy analizie poddano kilka wybranych architektur sieci neuronowych zastosowanych w agentach uczonych metodami A2C oraz PPO. Zastosowanie zaawansowanych modeli głębokiego uczenia jest uzasadnione specyfiką problemu, która została omówiona w poprzednich rozdziałach: agent działający w środowisku Blood Bowl musi przetwarzać złożone, bogate w informacje stany gry, składające się zarówno z danych przestrzennych (układ zawodników na planszy, strefy zagrożeń, pozycja piłki), jak i danych nieprzestrzennych (liczba tur, dostępne przerzuty, aktywne procedury gry). Złożoność decyzyjna tej gry – obejmująca zmienne drzewo decyzyjne, ogromny branching factor oraz liczne zależności wieloetapowe – wymaga zastosowania modeli zdolnych do skutecznego wydobywania informacji z reprezentacji wielowymiarowych.

W tym kontekście sieci konwolucyjne (CNN) stanowią naturalny wybór jako moduły ekstrakcji cech przestrzennych. Umożliwiają one analizę map cech o wysokiej rozdzielczości i wykrywanie lokalnych oraz globalnych wzorców. W celu zwiększenia efektywności przetwarzania przestrzennego oraz stabilności uczenia, w pracy zastosowano trzy odmienne podejścia architektoniczne:

\begin{itemize}
	\item ResNet-CNN – klasyczna architektura rezydualna oparta na blokach ResNet o stałej liczbie kanałów, znana z doskonałej stabilności i efektywnej reprezentacji w zadaniach wizji komputerowej.
	\item CNN z blokiem Inception – upraszczając, równoległe operacje konwolucyjne o różnych rozmiarach filtrów umożliwiają analizę cech w wielu skalach jednocześnie (inspiracja GoogLeNet).
	\item Inception z połączeniami rezydualnymi – dodanie tzw. skip connections ułatwia propagację gradientu i pozwala stabilnie trenować głębsze sieci.
\end{itemize}

\section{Ogólna struktura agentów A2C i PPO}

\subsection*{Polityka i wartość (Actor–Critic)}
Zarówno A2C, jak i PPO są algorytmami typu actor–critic, co oznacza, że model sieci neuronowej posiada dwie odrębne głowy wyjściowe:

\begin{itemize}
	\item politykę (actor), która generuje rozkład prawdopodobieństwa akcji,
	\item wartość (critic), która ocenia jakość aktualnego stanu.
\end{itemize}

Formalnie polityka jest funkcją:
\[\pi_\theta(a|s)\]
gdzie $\theta$ oznacza parametry sieci, $s$ jest stanem, a $a$ - akcją. Krytyk aproksymuje wartość stanu:
\[V_\theta(s)\]
czyli oczekiwaną zdyskontowaną sumę przyszłych nagród przy stosowaniu polityki $\pi_\theta$. Takie rozdzielenia funkcji decyzyjnej i oceniającej poprawia stabilność uczenia, ponieważ błędy krytyka dostarczają dodatkowego sygnału uczącego dla aktora.

\subsection*{Współdzielenie części konwolucyjnej}
